% Example usage of mko3.cls — compile with XeLaTeX
% Requires: TH Sarabun New font installed
\documentclass{mko3}

\begin{document}

\mkothreetitle

% ----- 1. Course info -----
\coursecode{XX12345}
\coursename{ชื่อรายวิชาตัวอย่าง}
\program{หลักสูตรวิศวกรรมศาสตรมหาบัณฑิต สาขาวิศวกรรมคอมพิวเตอร์}
\coursegroup{วิชาเฉพาะด้าน}
\semester{2 / 2568}
\instructor{ผศ.ดร.ตัวอย่าง นามสมมุติ}
\contact{example@uru.ac.th \quad โทร. 0-XXXX-XXXX}
\coursedesc{คำอธิบายรายวิชาเขียนรายละเอียดเนื้อหาวิชา จุดเน้น
ทักษะที่ผู้เรียนจะได้รับ และความเชื่อมโยงกับรายวิชาอื่น ๆ
ในหลักสูตร}
\courseinfotable

% ----- 2. CLOs domain/level -----
\begin{clodomaintable}
  \cloitem{CLO1}{อธิบายแนวคิดและทฤษฎีพื้นฐานของรายวิชาได้}{Cognitive}{3}
  \cloitem{CLO2}{ออกแบบและพัฒนาระบบจากโจทย์ที่กำหนดได้}{Psychomotor}{4}
\end{clodomaintable}

% ----- 3.1 CLO ↔ PLO mapping -----
\begin{cloplotable}{3}{ & \textbf{PLO1} & \textbf{PLO2} & \textbf{PLO3}}
  \textbf{CLO1} & X &   & X \\ \hline
  \textbf{CLO2} &   & X & X \\ \hline
\end{cloplotable}

% ----- 3.2 Bloom + LLL -----
\begin{bloomtable}
  \bloomrow{CLO1}{3}{2}{}{1}
  \bloomrow{CLO2}{4}{}{3}{2}
\end{bloomtable}

% ----- 4. Content analysis -----
\begin{contentanalysistable}
  \contentrow{หลักการพื้นฐาน}{Cognitive 2}{อธิบาย}{ในบริบทของระบบจริง}{3}
  \contentrow{การออกแบบ}{Psychomotor 3}{ออกแบบ}{ตามโจทย์ที่กำหนด}{6}
\end{contentanalysistable}

% ----- 5. Teaching & assessment -----
\begin{teachingtable}
  \teachrow{CLO1}{บรรยาย, อภิปราย, กรณีศึกษา}{แบบฝึกหัด, ข้อสอบ}
  \teachrow{CLO2}{ปฏิบัติการ, โครงงาน}{การประเมินโครงงาน, การนำเสนอ}
\end{teachingtable}

% ----- 6. Assessment plan -----
\begin{assessmenttable}
  \assessrow{สอบก่อนเรียน (ถ้ามี)}{--}{0}
  \assessrow{สอบระหว่างเรียน}{CLO1}{30}
  \assessrow{สอบปลายภาค}{CLO1, CLO2}{40}
  \assessrow{กิจกรรมการอื่น ๆ (ถ้ามี)}{CLO2}{30}
\end{assessmenttable}

% ----- 7. Support resources -----
\begin{supporttable}
  \supportrow{ห้องเรียน}{1 ห้อง}{พร้อมใช้}{ปานกลาง}
  \supportrow{ห้องสำหรับการเรียนรู้}{1 ห้อง}{พร้อมใช้}{ทันสมัย}
  \supportrow{อุปกรณ์}{ครบ}{พร้อมใช้}{ทันสมัย}
  \supportrow{เครื่องมืออื่น ๆ เช่น เอกสาร ตำรา วารสาร ที่แนะนำ}{ครบ}{พร้อมใช้}{ทันสมัย}
\end{supporttable}

% ----- Signatures -----
\signatureblock{ผศ.ดร.ตัวอย่าง นามสมมุติ}{ผศ.ดร.รองคณบดี ฝ่ายวิชาการ}

\end{document}
